\section{One Euler factor and five kinds of structure}
\label{sec:opening}

For a prime number $p$, put
\[
 E_p(s)=\frac1{1-p^{-s}},\qquad
 p^{-s}=\exp(-s\log p),\qquad \operatorname{Re}s>0,
\]
where $\log p$ is real. The geometric series defining $E_p$ counts
arbitrary occupation of one mode of energy $\log p$. It also occurs in
integration over the rational local field. The question is what, beyond
this scalar value, either construction determines. An arithmetic
cohomological explanation of the zeta function must retain operations
that the value alone does not record: multiplication, local restriction,
duality, and the choice of a state.

The local integral has complex coefficients and the $p$-adic field
$\mathbb Q_p$ as its domain. Its test space consists of locally constant,
compactly supported functions, with the inductive-limit topology of
finite-dimensional spaces of functions of bounded support and fixed
local constancy. These functions have degree zero and pointwise product.
Let $\mathbb Z_p$ be the valuation ring, take $|p|_p=p^{-1}$, and
normalize multiplicative Haar measure $d^\times x$ by
$\operatorname{vol}(\mathbb Z_p^\times)=1$. The chosen test is
$\one_{\mathbb Z_p}$. Its support and this measure are essential:
changing to additive measure changes the normalization.

The occupation space is the complex Hilbert space
$\mathscr H_p=\ell^2(\mathbb N_0)$ with orthonormal vectors $e_r$.
The number operator in energy units is
\[
 N_pe_r=r\log p\,e_r,\qquad
 \mathcal D(N_p)=
 \left\{a:\sum_{r\ge0}r^2(\log p)^2|a_r|^2<\infty\right\}.
\]
Here $r$ is occupation number, not cohomological degree. The finite
sequences form a common algebraic domain. Bounded operators multiply by
composition; the Hilbert space itself has no specified multiplication.
The trace of $e^{-sN_p}$ exists exactly for $\operatorname{Re}s>0$.
The basis vector $e_0$ is the vacuum. At real inverse temperature
$\beta>0$, the equilibrium density is the normalized positive operator
$(1-p^{-\beta})e^{-\beta N_p}$.

A third carrier is the one-dimensional complex Hilbert space
$V_p=\mathbb C e_p$, with bounded operator $K_p(s)e_p=p^{-s}e_p$.
It is concentrated in degree zero. Its ordinary determinant, normalized
to one as real $s\to+\infty$, is $\det(1-K_p(s))=1-p^{-s}$.
Thus the comparison of these three constructions is the scalar identity
\begin{equation}
 \int_{\mathbb Q_p^\times}\one_{\mathbb Z_p}(x)|x|_p^s\,d^\times x
 =\operatorname{Tr}(e^{-sN_p})
 =\det(1-K_p(s))^{-1}=E_p(s).
 \label{eq:opening-euler}
\end{equation}
The first equality is proved with the local-field and measure
constructions in Proposition~\ref{prop:local-tate}; the occupation sum
is the same convergent geometric series. No arrow between algebras or
complexes has been inferred from this equality.

Multiplication enters separately through effective divisors. Over a
commutative coefficient ring $R$, let $A_p=R[u_p]$, with unit $1$,
augmentation $\epsilon(u_p)=0$, discrete coefficients and algebraic
tensor products. Its elements have homological degree zero and the
weight grading $\operatorname{wt}(u_p^r)=r$; multiplication adds weights.
The bar construction is a complex built from this multiplication and
augmentation. Over $\mathbb C$, completing the monomials to an
orthonormal basis gives the occupation space above, and multiplication
by $u_p$ gives the unilateral shift. The square-zero algebra
$\mathbb C\oplus\bigoplus_{r>0}\mathbb C e_r$, with the same unit,
augmentation, weights and energy operator, has the same heat trace
but a different bar homology (Proposition~\ref{prop:two-bars}).
Consequently the scalar factor does not determine the multiplication
needed for arithmetic factorization.

An arithmetic cohomology class has a fifth carrier. For a totally
imaginary number field $F$, set $X=\operatorname{Spec}\mathcal O_F$.
For a positive integer $n$, assume $F$ contains a chosen primitive
$n$th root of unity. Ordinary sheaf cohomology on the small \'etale
site of $X$ has finite discrete coefficient ring $\mathbb Z/n$,
cohomological grading and cup product. A finite-group torsor on $X$
pulls a group-cohomology class back to degree three; the chosen root
normalizes its trace in $\tfrac1n\mathbb Z/\mathbb Z$. Removing a
nonempty finite set of places containing those above $n$ gives local
cochains and finite gluing maps. These are arithmetic operations, but
their coefficients and grading differ from those of the complex
Euler factor. A determinant comparison would require a realization
functor, an operator and a trace-compatible map; none is specified by
the existence of the degree-three class.

The analytic comparison can nevertheless be made exact. On the Hilbert
space with one basis vector for each prime, the diagonal $K(s)$ has
singular values $p^{-\operatorname{Re}s}$. Its $m$th regularized
determinant is defined precisely when $m\operatorname{Re}s>1$.
Analytic continuation preserves the normalized germ but may make it
multivalued. The logarithmic derivative defines a period character;
its kernel determines the smallest pointed surface on which the germ
is single-valued. For several orders the character is vector-valued.
On $\operatorname{Re}s>1/5$, orders two and three have individual
degrees six and twelve, but their joint surface has degree twenty-four.
This finite example distinguishes simultaneous continuation from
replacing several degree conditions by their least common multiple.

Write $\zeta(s)=\sum_{n\ge1}n^{-s}$ on $\operatorname{Re}s>1$
and $\GammaR(s)=\pi^{-s/2}\Gamma(s/2)$, where
$\Gamma(z)=\int_0^\infty e^{-u}u^{z-1}\,du$ for $\operatorname{Re}z>0$.
The real Gaussian factor has a different operator explanation.
An equally spaced additive spectrum gives a reciprocal Gamma
determinant, with spacing and scale determining its constant. At
spacing $2$ and scale $2\pi$ the value is $\sqrt2/\GammaR(s)$.
Together with the order-one prime determinant $1/\zeta(s)$ this
recovers $2^{-1/2}\GammaR(s)\zeta(s)$. This is an analytic identity
on specified carriers. Deninger's stronger question asks for
arithmetically defined cohomology and an operator whose alternating
determinant has this value. Its degree-one determinant is determined
conditionally in Theorem~\ref{thm:deninger}; the theorem does not
construct that operator or its eigenvalue multiplicities.

Two positive constructions clarify what an arithmetic comparison
could preserve. On a smooth projective curve over a finite field,
Hall multiplication and its adjoint give an operator exchange with
the curve zeta function as coefficient, before any trace is taken.
For the Riemann zeta function, the complete Weil functional has
infinite-dimensional positive subspaces on sufficiently narrow finite
unions of intervals. These subspaces retain all polar, prime-power
and archimedean terms. Their supports are not stable under convolution,
so they do not supply the full positive representation required by
Weil's criterion. Arithmetic twisting, full positivity and
cohomological spectral realization are therefore distinct targets,
with the comparison maps stated where their carriers are constructed.

The inner-product convention is part of every operator comparison.
The occupation and oscillator spaces use an inner product linear in
the first argument. The Hall pairing and the positive Weil pairing use
one linear in the second.
If $\langle x,y\rangle_1$ is the first convention, put
$\langle x,y\rangle_2=\langle y,x\rangle_1$ to pass to the second.
The adjoint operator is unchanged: conjugating
$\langle Ax,y\rangle_1=\langle x,A^*y\rangle_1$ gives the
corresponding identity in the second convention. Rank-one formulas
must also be transported: $h\mapsto\langle h,g\rangle_1 f$ becomes
$h\mapsto\langle g,h\rangle_2 f$. This fixes the convention before
the equilibrium, Hall and finite-rank trace calculations are compared.
